\section{Hedges and domain-explicit completeness: full proofs}
\label{sec:annex-p16-hedges-and-completeness}

\begin{annexitem}{annex:p16.hedge-singleton-reduction}{thm:p16-hedge-singleton-reduction}{Exact removal and restoration of constant coordinates}
\begin{proof}
Every total equation with codomain \(S_i=\{c_i\}\) returns
\(c_i\) for every argument. In each remaining equation substitute
\(c_i\) for every parent in \(C\). Retain all original latent
inputs and intersect their child sets with \(N_*\). Their joint
law remains a product. Directed parents and shared latent inputs
between retained vertices give exactly \(G[N_*]\); no path
through a constant coordinate creates a new parent or a new common
input. The resulting functions are total on their reduced domains.

For each fixed latent tuple, evaluate both models in a topological
order. Constant coordinates agree with their unique values; induction
then makes every retained coordinate agree. The same induction applies
after an arbitrary intervention: a legal action on a constant coordinate
replaces a constantly valued equation by that same value, and every
other replaced equation has the same assigned value in both models.
Thus all observed and interventional laws agree after deleting the
constant coordinates. On a response set containing constants, the
omitted coordinates have their unique assignment with probability one.

Conversely take any reduced model. Add the constant coordinates and
their constant equations, and add all original declared directed edges
incident with them as ignored arguments. For every missing declared
bidirected edge incident with a constant coordinate, add one independent
constant latent input with exactly that pair of children; both equations
ignore it. Keep the reduced model's other latent inputs. This gives
exactly the required full declared mixed graph, without changing the
law of any retained coordinate. Constants could also be fair ignored
inputs; no particular latent cardinalities were prescribed. This is
why the statement uses the graph-defined class, rather than a class
with separately fixed latent incidences or cardinalities.

There is a bijection between full observed assignments and reduced
assignments, obtained by adjoining \(c\). Their masses are equal,
so positivity is equivalent. If the reduced kernel is identified,
apply reduction to any two original models with the same observed
law to obtain equal reduced effects, then restore the constant response
coordinates. If the original kernel is identified, extend any two
reduced models with the same observed law by the construction just
given and project the resulting equal effects back. This proves both
directions of identification equivalence, including empty reduced
action or response sets.
\end{proof}
\end{annexitem}

\begin{annexitem}{annex:p16.hedge-failure-extraction}{thm:p16-hedge-failure-extraction}{The trace normal form and a common forest construction}
\begin{proof}
Work after singleton elimination and use the reduced initial sets
\(N,X_0,Y_0\). A failed run has a nonempty response, since an
empty response reduces to the empty graph and returns one. First
let \(R_0=\operatorname{an}_G(Y_0)\). Branch 2 either makes
this restriction or finds it already satisfied. Its local actions
are \(X_0\cap R_0\). Every directed path ending in \(Y_0\)
lies in \(R_0\), so
\[
 \operatorname{an}_{G[R_0\setminus X_0]}(Y_0)
   =\operatorname{an}_{G[N\setminus X_0]}(Y_0)=A.
\]
Branch 3 therefore makes the nonaction set precisely \(A\),
if it was not already \(A\), keeping the same random set and
kernel. Its new actions outside \(X_0\) are visible in the frame.
Every vertex of \(A\) has a directed path to \(Y_0\) inside
\(A\), so a subsequent ancestral reduction that retains \(A\)
in its random graph cannot discard any vertex of \(A\).

If \(G[A]\) has at least two districts, branch 4 now splits them.
The selected child on a failed path has response \(D\) and
nonaction set \(D\), for a maximal district \(D\) of \(G[A]\).
Ancestral restriction always keeps the entire response \(D\),
and branch 7 always keeps the entire nonaction set \(D\).
Thus in this child and every descendant, the nonaction and response
sets remain \(D\). Branch 3 cannot add any action because every
vertex of the nonaction set is itself a response. Branch 4 cannot
split because the same induced graph \(G[D]\) remains connected
by bidirected edges. The only recursive edges thereafter are 2 and 7.

If instead \(G[A]\) has one district, it never splits later:
branch 7 preserves the whole nonaction set \(A\), and a later
branch 2 keeps it by its directed paths to \(Y_0\) inside
\(A\). Branch 3 has no new nonancestor to add. Subsequent
recursive edges are again only 2 and 7, with response \(Y_0\).
This proves the two claimed trace forms. In particular, neither
the enlargement of actions in branch 3 nor the changed response
in branch 4 is described as a subset of the original input set.

At failure the already proved ID invariant gives
\[
 S\subsetneq R,\qquad R=\operatorname{an}_{G[R]}(K),\qquad
 S=\operatorname{an}_{G[S]}(K),
\]
and bidirected connectivity of both induced graphs. Suppose
\(R\cap X_0=\varnothing\). Every vertex in \(R\) then has
a path to \(K\subseteq A\), staying outside \(X_0\), and
every member of \(A\) has a path to \(Y_0\) outside \(X_0\).
Concatenating directed paths, or deleting repetitions, puts \(R\)
inside \(A\). In the unsplit case this says \(R\subseteq S\),
a contradiction. In the split case \(R\) is bidirected-connected
inside \(A\) and contains the maximal district \(D=S\);
maximality forces \(R\subseteq D\), again a contradiction.
Thus \(R\cap X_0\ne\varnothing\), while \(S\cap X_0=\varnothing\)
because \(S\subseteq A\subseteq N\setminus X_0\).

For every \(s\in S\setminus K\), choose the first edge of
a shortest directed path from \(s\) to \(K\) within \(S\).
For every \(r\in R\setminus S\), choose the first edge of
a shortest directed path from \(r\) to \(S\) within \(R\).
Give vertices of \(K\) no selected child. These choices exist
by the displayed ancestor identities. Distances decrease until the
path enters \(S\), then decrease toward \(K\). Thus the
selected directed graph has at most one child per vertex, exactly
the sinks \(K\), and keeps \(S\) closed under its descendants.
It is acyclic as a subgraph of \(G\).

Choose a bidirected spanning tree of \(G[S]\). Treat its vertices
as an initial connected set in \(G[R]\). While a vertex of
\(R\) is missing, connectivity gives an edge from the present
set to a missing vertex along a path to it; add that edge and
vertex. Each addition preserves a tree and eventually covers \(R\).
There are no new edges with both endpoints in \(S\), so the
restriction to \(S\) is the originally chosen tree. Together
with the selected directed edges these trees give the nested
\(K\)-rooted C-forests \(F,F'\).

Finally \(K\subseteq A=\operatorname{an}_{G[N\setminus X_0]}(Y_0)\).
All four hedge conditions hold for the original query. This is a
constructive version of the source Theorem 6 route. It selects
common roots explicitly instead of assuming that arbitrary directed
pruning preserves identical root sets in the two graphs.
\end{proof}
\end{annexitem}

\begin{annexitem}{annex:p16.hedge-positive-parity}{thm:p16-hedge-positive-parity}{The complete parity construction and its positive smoothing}
\begin{proof}
First note that \(S\) is descendant-closed in the selected directed
forest on \(T\). Indeed every \(s\in S\setminus K\) has
a child in \(F'\), which must be its unique child in \(F\),
and every \(k\in K\) has no child in either. Thus no selected
directed arrow leaves \(S\). Every vertex of either forest has
one selected path to its root in \(K\).

We require an elementary fact about a tree on \(m\) vertices.
For \(m>1\), an endpoint of a longest simple path has degree
one: another incident edge would either extend the path or form a
cycle. Removing that endpoint and its edge leaves a tree, so induction
also proves that a tree on \(m\) vertices has \(m-1\) edges.
Its edge bits map to vertex bits by assigning to each vertex the
parity of its incident edges. The resulting vertex vector has even
total parity because each edge is counted twice. The map is
injective: if the incidence parity is zero at every vertex, an
end vertex forces its sole incident edge to zero; remove that edge
and vertex and repeat until the tree is exhausted. There are
\(2^{m-1}\) edge assignments and \(2^{m-1}\) even vertex
vectors, so the map is a bijection onto the latter set. The case
\(m=1\) has the single empty edge assignment and vertex value zero.

In model 1 define, for a proposed observed vector,
\[
 z_i=v_i\oplus\bigoplus_{j\text{ selected parent of }i}v_j.
\]
The map \(v\mapsto z\) is bijective: solve its coordinates
in topological order. Since every nonroot has exactly one child,
\(\bigoplus_i z_i=\bigoplus_{k\in K}v_k\).
Conditional on \(E=e\), the vector \(z\) is the incidence
vector of the fair edge bits with \(e\) added at \(k_0\).
The tree fact shows that it is uniform over vectors of total parity
\(e\). Hence \(v\) is uniform over all \(2^{|T|-1}\)
assignments whose root parity equals \(e\).

For model 2, put \(O=T\setminus S\). There are \(|S|-1\)
tree edges within \(S\) and \(|O|\) remaining tree edges.
Contract the connected subtree on \(S\) to one vertex. The
remaining edges form a tree on \(O\) and that contracted vertex.
Its edge-incidence map restricted to coordinates \(O\) is a
bijection onto all bit vectors on \(O\). To see injectivity,
zero incidence on \(O\) forces zero incidence at the contracted
vertex by even total parity, after which the tree fact forces
every edge bit to zero. Domain and codomain each have \(2^{|O|}\)
elements, proving bijectivity.

The equations on \(O\) have no selected directed parents in
\(S\), by descendant closure. Topological solution transforms
their uniform incidence vector bijectively, so \(V_O\) is uniform
on all \(2^{|O|}\) assignments. It uses only the edge bits
outside the \(S\) subtree. The equations on \(S\) in model 2
use only its subtree edge bits and \(E\). By the preceding
argument within \(S\), they are conditionally uniform on root
parity \(e\). The two latent groups are independent. Multiplying
their masses gives the same conditional uniform law on \(T\)
as in model 1. Mixing \(e=0,1\) with weights \(1-\epsilon,\epsilon\)
proves the displayed common full law, which is strictly positive.

Now fix a nonempty \(Z\subseteq O\) to any bit assignment.
Let \(B\) be the vertices whose unique selected path to \(K\)
meets \(Z\), including \(Z\) itself. Its complement \(T\setminus B\)
contains all of \(S\), so both sets are nonempty. On expanding
the model-1 equations after the intervention, an incident edge bit
contributes to root parity once for each endpoint whose path avoids
\(Z\). Its coefficient is therefore one exactly when its tree
edge crosses between \(B\) and \(T\setminus B\). At least
one crossing edge exists because the selected bidirected tree is
connected. That edge contributes an independent fair bit with
coefficient one. Fixed action values contribute a constant parity,
and \(E\) contributes once because \(k_0\notin Z\).
Regardless of all other terms, the resulting root parity is fair.

In model 2 the whole submodel on \(S\) ignores every input
outside \(S\). Its equations are unchanged by actions on \(Z\).
The same incidence cancellation gives root parity exactly \(E\).
The two probabilities are thus \(1/2\) and \(\epsilon\).

This retains the source Theorem 4 parity route while proving equality
of the \emph{full} observed laws, including the vertices outside
the smaller forest. Setting \(\epsilon=0\) recovers its uniform
even-root-parity construction. The positive theorem instead uses
\(0<\epsilon<1/2\); no zero mass is offered as positive support.
\end{proof}
\end{annexitem}

\begin{annexitem}{annex:p16.hedge-alphabet-lift}{thm:p16-hedge-alphabet-lift}{A decoded structural model on prescribed finite alphabets}
\begin{proof}
For each \(i\) and \(b\in B_i\), let \(R_{i,b}\) be a
private random choice from the nonempty fiber \(d_i^{-1}(b)\),
with mass function \(w_i\) on that fiber. Take all these choices
mutually independent and independent of the original latent inputs.
Their finite tuple \(R_i=(R_{i,b})_{b\in B_i}\) can be one
private latent input at \(i\). Denote the original observed variable
with alphabet \(B_i\) by \(W_i\). If its equation is
\(W_i=f_i(W_{\operatorname{pa}(i)},U_{\Lambda_i})\), define
the lifted total equation by
\[
 V_i=R_{i,\,f_i(d(V_{\operatorname{pa}(i)}),U_{\Lambda_i})}.
\]
Applying \(d_i\) to its output returns the indicated value of
\(f_i\). Induction in a topological order shows that decoded
observed variables solve exactly the original equations on the same
original latent tuple. Extra private choices do not change the
declared directed or bidirected edges.

For a full proposed output vector \(v\), the original latent
constraints are exactly those for decoded vector \(d(v)\).
Conditional on these constraints the independent private choices
contribute \(\prod_i w_i(v_i)\). Summing over the original
latent tuple gives the displayed formula for \(\widetilde p\).
Both original models have the same \(p\), and every factor is
positive. Summing all lifted assignments gives one because each
fiber's weights sum to one and the decoded law is normalized.

After replacing the equations on \(X\) by \(x\), their decoded
coordinates are fixed to \(d(x)\). Induction on the remaining
vertices again gives the original intervened solution after decoding.
Private choices on the action vertices are unused. For a response
assignment \(v_Y\), summation over all other random coordinates
and their fibers therefore gives
\[
 \widetilde P_x(v_Y)
   =P_{d(x)}(d(v_Y))\prod_{i\in Y}w_i(v_i).
\]
For a decoded response event, sum this formula over its fibers;
each product of fiber weights sums to one, proving exact preservation
of its probability. A different original probability remains different.
Every required latent set is finite, and all new roots are private,
so the construction stays in the declared model class.
\end{proof}
\end{annexitem}

\begin{annexitem}{annex:p16.hedge-binary-routing}{thm:p16-hedge-binary-routing}{Embedding a hedge and transporting its parity along external paths}
\begin{proof}
Start with the two positive forest models on \(T\). For every
\(h\in H\) use the specified parity equation and private bit.
Give each remaining observed vertex outside \(T\cup H\) an
independent fair private bit, ignoring its other arguments. Declare
all directed edges of the original graph. For each bidirected edge
not yet supplied by the selected forest roots, add an independent
constant latent input shared by just its two endpoints, ignored by
both equations. These additions supply exactly the original declared
mixed graph and introduce no unwanted edge. All latent inputs remain
independent. In particular an intervention elsewhere cannot affect
the forest through a declared argument its equations ignore.

No routing edge enters \(K\), and \(H\) lies outside \(T\).
Thus the equations and joint law on \(T\) stay unchanged. Given
the whole assignment on \(T\), the routing variables have the
same product of conditional noise masses in both models, evaluated
in topological order. Each of those masses is either \(\epsilon_h\)
or \(1-\epsilon_h\), both positive. The remaining coordinates
contribute independent factors \(1/2\). It follows that the
full observed laws are equal and positive at every assignment.

The action set \(X\) misses the routing vertices and all of
\(S\), but meets \(T\setminus S\). Its actions within
\(T\) are precisely a nonempty set \(Z\) covered by the
positive parity theorem. Actions outside \(T\) replace only
the irrelevant independent equations, since the routing forest
avoids \(X\). Sum the routing equations modulo two. Each
nonterminal routing value occurs once as an output and once at
its unique child's input. Cancellation yields
\[
 \bigoplus_{j\in J}V_j
   =\bigoplus_{k\in K}V_k\ \oplus\!
      \bigoplus_{h\in H}E_h.
\]
This identity includes roots already in \(J\) by their zero-length
paths. The noise bits are independent of the forest variables and
of each other, before and after the intervention. Since
\(\mathbb E[(-1)^{E_h}]=1-2\epsilon_h\), the expected sign
of the displayed parity is the product of the root-parity sign
expectation and those noise expectations. The first root expectation
is zero; the second is \(1-2\epsilon\). Their discrepancy is
strictly positive and is visible on \(J\subseteq Y\).
This proves binary nonidentification at the original query, not
only at its failed child. Apply the alphabet lift to extend the
same witness to any prescribed value sets of size at least two.
\end{proof}
\end{annexitem}

\begin{annexitem}{annex:p16.hedge-four-state-routing}{thm:p16-hedge-four-state-routing}{Separate core and relay coordinates on a fixed four-value domain}
\begin{proof}
First take exactly four values at every observed vertex, identified
with \(\{0,1\}^2\). The observed variable at \(i\) is a
single pair \((C_i,L_i)\). This does not add observed vertices.
The two core models on \(T\) are the positive forest models.
For \(i\notin T\), take \(C_i\) to be an independent fair
private bit. Core equations ignore all relay coordinates and all
directed edges not selected by the original forest.

Let \(A=\operatorname{an}_{G[N\setminus X]}(Y)\). For every
\(i\in A\setminus Y\), select the first edge on a shortest
directed path from \(i\) to \(Y\) in \(G[N\setminus X]\).
Give vertices in \(Y\) no selected outgoing edge. Distances
decrease, so every selected path ends in \(Y\). Let \(H\)
be the union of the selected paths starting at \(K\), including
their endpoints, and let \(J\subseteq Y\) be the sinks reached.
The selected graph on \(H\) has at most one child per vertex
and avoids \(X\). Paths can merge and can pass through other
vertices of \(K\), through vertices of \(T\), or through
vertices outside \(T\). None of these possibilities is excluded.

For every \(i\in H\), take an independent private bit \(E_i\),
with \(0<\eta_i=\mathbb P(E_i=1)<1/2\), independent of all
core latent inputs. Define
\[
 L_i=\left(\bigoplus_{j\text{ selected routing parent of }i}L_j\right)
       \oplus \bigl(\mathbf1_{\{i\in K\}}C_i\bigr)\oplus E_i.
\]
At vertices outside \(H\), use an independent fair private bit
as \(L_i\). The product \(\mathbf1_{\{i\in K\}}C_i\)
here is multiplication of bits, so it inserts the core bit exactly
at the vertices in \(K\).

These are legal total structural equations for the observed pairs.
At vertex \(i\), its core value is first computed from its
observed parents' core coordinates and its latent inputs; its relay
value is then computed using that already computed bit and the
parents' relay coordinates. This is a single function of the original
observed parents and latent inputs. It creates no self-parent edge
and no extra observed vertex. All selected parent edges belong to
\(G\), and all added noise inputs are private. Supply remaining
declared edges by ignored arguments and pair-specific ignored inputs
as in the external-routing proof.

The full core law is the common positive law on \(T\) times
independent fair masses outside \(T\), hence is equal in the
two models and positive everywhere. Conditional on any full core
assignment \(c\), the independent relay noise variables retain
their product law. In topological order the conditional mass of a
proposed relay vector \(l\) is the product of the appropriate
\(\eta_i\) or \(1-\eta_i\) at every \(i\in H\) and
the factors \(1/2\) elsewhere. The same product occurs in both
models and is positive for every \(l\). Thus the full observed
\emph{pair} law is equal and strictly positive on every four-valued
assignment, including all combinations of core and relay coordinates.

Now intervene at the original set \(X\), fixing the whole pairs
to arbitrary specified values. The resulting core equations are
the original core models with their core coordinates fixed at
\(X\). Their action set in \(T\) is nonempty and disjoint
from \(S\). The root core parity therefore has expected sign
zero in model 1 and \(1-2\epsilon\) in model 2. Every routing
equation is still present because \(H\cap X=\varnothing\).
Sum those equations modulo two. Each relay value at a nonterminal
vertex occurs twice, once as output and once at its unique child;
each root core bit is injected exactly once, even when a path from
one root passes through another. Hence
\[
 \bigoplus_{j\in J}L_j
   =\bigoplus_{k\in K}C_k\ \oplus\!
     \bigoplus_{i\in H}E_i.
\]
All relay noises remain independent of the core model after the
intervention. The expected signs of this parity are respectively
zero and
\((1-2\epsilon)\prod_{i\in H}(1-2\eta_i)>0\).
It is an event measurable from the observed pairs on
\(J\subseteq Y\), proving different original response laws.

It remains to justify the stronger local cardinality assertion.
Put \(I=(T\setminus K)\cap H\). Retain both coordinates
only at \(I\). Retain just the core coordinate on \(T\setminus H\),
just the relay coordinate on \((H\setminus T)\cup K\), and
one independent core coordinate outside \(T\cup H\).
Every selected core parent belongs to \(T\setminus K\),
because the roots \(K\) have no selected core child. Its core
coordinate is therefore still available, either alone or in a pair.
Every selected relay parent belongs to \(H\), where the relay
coordinate is retained. At \(k\in K\), compute its core
function from these available parent core coordinates and original
latent inputs, substitute that expression into its relay equation,
and output only the relay value. Thus every retained output remains
a total function of its declared observed parents and latent inputs;
the omitted core at a root is an internal calculation, not an
extra hidden endogenous coordinate assumed available to other nodes.

These equations produce precisely the coordinate projection of
the four-valued model, as topological induction verifies. The
projected observed laws are equal and strictly positive at every
assignment: each such mass is a nonempty sum of positive masses
in the four-valued law. The original action set misses \(H\),
so no action fixes an omitted root core or a retained relay. At
action vertices in \(T\setminus H\), fixing their retained
core has the same effect as fixing the full old pair to that core
and any relay value, since that relay was unused. Outside \(T\cup H\)
the extra coordinate was also unused. Consequently the intervened
projected laws agree with projections of the corresponding four-valued
intervention laws. The discrepant parity on \(J\) uses relay
coordinates retained everywhere on \(J\subseteq H\), so its
discrepancy survives. This proves a witness with four-element alphabets
at \(I\) and binary alphabets elsewhere.

For prescribed domains meeting those respective lower bounds, choose
surjections onto these two- or four-element alphabets and apply the
proved alphabet lift. Equality, strict positivity and the decoded
intervention discrepancy all survive. In particular domains of size
at least four everywhere satisfy the requirements.
The construction is a full routing proof for the source Theorem 4
obstruction on this stated domain. It makes explicit the additional
coordinate used when routes reenter the original forest, so it
does not present a four-valued witness as a fixed-binary witness.
\end{proof}
\end{annexitem}

\begin{annexitem}{annex:p16.hedge-completeness}{thm:p16-hedge-completeness}{Propagation to the initial query and the resulting completeness classes}
\begin{proof}
The singleton theorem equates the original and reduced identification
problems and their positive observed-law classes. Work in the reduced
graph and retain the reduced original \(X,Y\). Termination and
the correctness of every successful expression have already been
proved in \cref{thm:p16-id-termination,thm:p16-id-success-soundness}.
Thus a successful run identifies the kernel in every class considered
here, independently of its value-set cardinalities.

For a failure, the extraction theorem produces forests forming a
hedge for the \emph{initial} \(X,Y\). Its trace proof establishes
both a genuine intersection with the initial action set and paths
from the common roots to the initial response that avoid those
actions. Thus subsequent model constructions can be made directly
on the original graph. One need not infer that a nonidentified
child factor remains nonidentified after arbitrary multiplication
or summation; such an inference alone would not be valid. The
constructed response event provides that missing implication.

If all reduced domains have size at least four, apply the four-state
routing theorem and the alphabet lift to this extracted hedge.
The resulting two models have the original declared graph and
the unchanged prescribed domains, strictly positive equal observed
laws, and different kernels for the original query. They therefore
belong to the exact class whose identification was in question.
The frame's added actions and reference values are not imposed on
these countermodels: the discrepancy holds under the original
actions themselves. Restoring singleton coordinates lifts both
models back to the original unreduced problem. This proves that
failure implies nonidentification in the first class. Together
with success soundness and termination it proves the equivalence.

For the full nonaction response \(Y=N\setminus X\), the common
root set of any extracted hedge lies in \(Y\), since the smaller
forest is disjoint from \(X\). It has the empty external routing
forest, with \(J=K\). The binary routing theorem therefore gives
a positive witness, and the binary alphabet lift permits all
remaining prescribed domains of size at least two. The same argument
proves failure nonidentification and hence completeness for the
second class. A specified empty full response has the unit kernel
and is handled by success before any failure arises.

For the third clause, an external routing forest supplies the premise
of the binary transport theorem. Alternatively, the stated local
cardinality condition supplies the premise of the refined routing
construction. Applying the appropriate theorem, the alphabet lift
and singleton restoration constructs the two original-query
countermodels. This is a sufficient failure
certificate on that domain, without claiming that the routing
condition is necessary or that every remaining failed trace has it.

Finally the complete source Lemma 4--8 derivations in
\cref{thm:p16-id-do-utility} turn every successful ID trace into
a finite derivation by the three sound do-calculus rules and ordinary
probability operations. All its conditional denominators are positive
by the proved execution invariant. In each complete class, an
identified kernel cannot have a failed run because that would
produce the countermodels just constructed. Its terminating run
must therefore succeed and supply such a derivation. This is the
source Theorem 7 inference, with the scope of the completeness
premise retained. Conversely the already proved soundness of the
rules makes a universal observed-law formula impossible whenever
the two covered countermodels disagree on its target. These are
unconditional intervention-kernel conclusions; no conditional-effect
or counterfactual identification algorithm is invoked.
\end{proof}
\end{annexitem}

\begin{annexitem}{annex:p16.hedge-all-effects}{thm:p16-hedge-all-effects}{The graph criterion for identifying all effects on nonconstant domains}
\begin{proof}
Work after the proved singleton reduction. Suppose no directed edge
has its endpoints in the same maximal district of the reduced
graph. If an ID run failed, its raw random graph on \(R\) would
be one district, and it would contain a vertex outside its nonaction
set with a directed path to its nonempty local response. At least
one directed edge of that path has both endpoints in \(R\).
They are connected by bidirected edges inside \(R\), hence lie
in the same maximal district of the reduced original graph. This
contradicts the hypothesis. Every run therefore terminates successfully.
Its expression and do-calculus derivation identify every query by
the previous success theorem.

Conversely suppose there is a directed edge \(a\to b\) whose
endpoints belong to one maximal district. Choose a simple bidirected
path from \(a\) to \(b\), with vertex set \(T\). Let
\(F\) contain only its bidirected path edges and the directed
edge \(a\to b\). Its roots are \(K=T\setminus\{a\}\).
Let \(F'\) delete \(a\) and its incident selected edges.
Because \(a\) is an endpoint of the simple bidirected path,
the remaining path is connected. There are no selected directed
edges in \(F'\), so its root set is the same \(K\).
These forests form a hedge for action \(\{a\}\) and response
\(N\setminus\{a\}\), since its roots are response vertices.
The positive binary witness with empty external routing, embedded
in the remaining independent coordinates and lifted to the prescribed
domains, disproves identification of this full nonaction kernel.
All observed domains have at least two values, as required for
that lift. Thus not every effect is identified. Restoring constants
proves the stated equivalence in the original graph-defined class.
\end{proof}
\end{annexitem}

\begin{annexitem}{annex:p16.hedge-worked}{thm:p16-hedge-worked}{Every branch and mass in the worked countermodels}
\begin{proof}
For the first query, the original random set is \(\{A,B,C,Z\}\),
the action is \(\{A\}\), and the response is \(\{Z\}\).
All vertices are directed ancestors of \(Z\), so branch 2
keeps them. After removing \(A\), the remaining graph has
\(B\to Z,C\to Z,B\leftrightarrow C\). All three remaining
vertices are response ancestors, so branch 3 adds nothing.
Branch 4 splits the two districts \(\{B,C\}\) and \(\{Z\}\).
The first child has response \(\{B,C\}\), actions \(\{A,Z\}\)
and the original kernel. Branch 2 restricts it to the ancestor set
\(\{A,B,C\}\), with kernel \(q_{ABC}=\sum_z p(a,b,c,z)\),
action \(\{A\}\), and a fixed \(Z\) argument which the kernel
ignores. Its remaining vertices \(B,C\) are one district and
its whole random graph is one district. It reaches branch 5 and
returns the raw pair on \(\{A,B,C\}\), \(\{B,C\}\).
This propagates through the saved ancestral and district-split
frames. Arithmetic is not applied to the failure object.

The chosen bidirected tree on \(\{B,C\}\) is its single edge.
Extending it adds \(A\leftrightarrow B\). The local response
is the whole smaller district, so both \(B,C\) are selected
roots and have no selected outgoing edge. On \(A\) choose
\(A\to B\), deleting the other allowed arrow \(A\to C\)
from this edge subgraph. This constructs the claimed nested forests,
which meet the \emph{original} action set at \(A\). The arrows
\(B\to Z,C\to Z\) give an external routing forest with
sinks \(\{Z\}\), so the binary propagation theorem applies.

In both models, observational substitution of \(A=U\) makes
\(B=V\oplus E\), while \(C=V\). Therefore \(A\) is an
independent fair bit and the root parity \(B\oplus C\) is
\(E\). For each triple \((a,b,c)\), its mass is \(3/16\)
when \(b=c\), and \(1/16\) when \(b\ne c\): the unique
required \(u,v\) contribute \(1/4\), and the required
\(E=b\oplus c\) contributes \(3/4\) or \(1/4\).
Conditional on that triple, \(Z=b\oplus c\oplus E_Z\).
For parity zero the two \(Z\) masses are therefore
\((3/16)(3/4)=9/64\) and \((3/16)(1/4)=3/64\).
For parity one they are \((1/16)(1/4)=1/64\) and
\((1/16)(3/4)=3/64\). Each parity class has four triples,
so the full sixteen masses sum to
\(4(9+3+1+3)/64=1\) and every cell is positive.

After \(A=a\), in model 1 the parity of \(B,C\) is
\(a\oplus U\oplus E\), which is fair because \(U\) is
fair and independent. Adding \(E_Z\) preserves fairness, so
\(\mathbb P^1_a(Z=1)=1/2\). In model 2 the parity remains
\(E\), giving
\[
 \mathbb P^2_a(Z=1)
   =\mathbb P(E\ne E_Z)
   =\tfrac14\tfrac34+\tfrac34\tfrac14=\tfrac38.
\]
This calculates the discrepancy in the original singleton response,
not merely the nonidentified child factor.

For the bow example, the binary common law is
\(p(a,0)=3/8\), \(p(a,1)=1/8\) for each binary \(a\).
The alphabet map sends 0 to binary 0 and sends 1,2 to binary 1;
the fiber weights are 1 at 0 and \(1/2\) each at 1,2.
The lifted treatment marginal is \((1/2,1/4,1/4)\), and the
lifted response marginal is \((3/4,1/8,1/8)\). Their product
is the displayed matrix: the first row is \((12,2,2)/32\),
and each remaining row is \((6,1,1)/32\). These entries are
positive and their total numerator is 32. Under every treatment
value, the decoded model-1 response is fair because its equation
contains the independent fair \(U\), whereas the decoded model-2
response is \(E\). The lifted event \(Y\in\{1,2\}\) has
the claimed probabilities by the alphabet theorem; each individual
nonzero response value has half the indicated event probability.

Finally, if \(A\) has just one legal value, its equation always
returns that value, whatever its declared parents. Its only legal
intervention does the same. Direct structural substitution gives
\(P_{\operatorname{do}(A=a)}(Y)=P(Y)\) in every such model.
The elimination step removes \(A\) and its incident edges and
runs ID with empty action set, returning that marginal. Running
the old unprocessed bow graph would still trigger its graph failure
test; the proved preprocessing prevents treating this as a valid
nonidentification certificate for a constant treatment.
\end{proof}
\end{annexitem}
